%%%%%%%%%%%%%%%%%%%%%%%%%%%%%%%%%%%%%%%%%%%%%%%%%%%%%%%
% % Lines starting with % are comments, which are ignored.
% % This is a handy way of indicating the date and version of
% % your document, to wit:
% %
% % LaTeX sample file
% % Modified March, 2002
% %
%%%%%%%%%%%%%%%%%%%%%%%%%%%%%%%%%%%%%%%%%%%%%%%%%%%%%%%
\DocumentMetadata{
  % PDF standards: ONE key whose value is a comma-separated list --
  % UA-2 (accessibility) plus A-4f (archival; the archival companion of
  % UA-2). Do NOT use separate pdfstandard=... keys: with ordinary
  % key-value parsing only the last one takes effect.
  pdfstandard={UA-2,A-4f},
  lang=en-US,
  tagging=on,
  tagging-setup={
    math/setup=mathml-SE
  }
}

\documentclass[12pt]{article}
% homework-macros.tex -- shared package setup for homework-template.tex
% (article class).
%
% NOTE ON PATHS: every template in this repository is stored one level up
% from where it is compiled (see README.md). The \input path below is
% resolved against the folder the document is compiled in, which sits
% one directory above this file.

% amsmath + amssymb + amsthm: full math support and proof environments
% (\newtheorem, \begin{proof}, ...).
\usepackage{amsmath, amssymb, amsthm}

% unicode-math: math typeset with OpenType math fonts (needed under
% XeLaTeX/LuaLaTeX so math matches the text fonts).
\usepackage{unicode-math}

% hyperref: clickable links (the template uses \href for the textbook
% problem sources); loaded after the math packages so it can patch them.
% Uses hyperref's default link colors.
\usepackage[colorlinks=true]{hyperref}

% Shared course math shorthands (\RR, \NN, \ZZ, \QQ, \grad, \curl, \tr,
% \der, \pder, ...) from the repository root.
\newcommand{\RR}{\mathbb{R}}
\newcommand{\NN}{\mathbb{N}}
\newcommand{\ZZ}{\mathbb{Z}}
\newcommand{\QQ}{\mathbb{Q}}
\newcommand{\vv}{\vec{v}}
\newcommand{\ww}{\vec{w}}
\newcommand{\xx}{\vec{x}}
\newcommand{\yy}{\vec{y}}
\newcommand{\zz}{\vec{z}}
\newcommand{\uu}{\vec{u}}
\newcommand{\basis}[1]{\vec{e_{#1}}}
\newcommand{\id}{\mathrm{id}}
\newcommand{\CC}{\mathbb{C}}
\renewcommand{\div}{\operatorname{div}}
\newcommand{\grad}{\operatorname{grad}}
\newcommand{\curl}{\operatorname{curl}}
\newcommand{\tr}{\operatorname{tr}}
\newcommand{\der}[2]{\dfrac{d #1}{d #2}}
\newcommand{\pder}[3][]{\dfrac{\partial^{#1} #2}{\partial #3}}


% % Title and author(s)
%%%%%%%%%%%%%%%%%%%%%%%%%%%%%%%%%%%%%%%%%%%%%%%%%%%%%%%
\title{Homework 4 - 09/17 Content}
\author{MATH 1271, Garrett Credi}
\date{09/17/2026}
%%%%%%%%%%%%%%%%%%%%%%%%%%%%%%%%%%%%%%%%%%%%%%%%%%%%%%%
% %
% % The next command allows your in import encapsulated
% % postscript files, .epsf or .eps files, which
% % contain vector graphic image data.
% %
%%%%%%%%%%%%%%%%%%%%%%%%%%%%%%%%%%%%%%%%%%%%%%%%%%%%%%%
% %   The first thanks indicates your affiliation
% %
% %  Just the name here.
% %
% % Your mailing address goes at the end.
% %
% % \thanks is also how you indicate grant support
% %
%%%%%%%%%%%%%%%%%%%%%%%%%%%%%%%%%%%%%%%%%%%%%%%%%%%%%%%

%%%%%%%%%%%%%%%%%%%%%%%%%%%%%%%%%%%%%%%%%%%%%%%%%%%%%%%
% $ The next command ensures the margins are 1 inch on
% % all sides.
% %
%%%%%%%%%%%%%%%%%%%%%%%%%%%%%%%%%%%%%%%%%%%%%%%%%%%%%%%
\begin{document}
\maketitle

Problems listed as AC are found in the \href{https://activecalculus.org/single1e/frontmatter.html}{Active Calculus} 
textbook at the end of the section listed.

Problems listed as APEX are found in the \href{https://opentext.uleth.ca/apex-calculus/apex-calculus.html}{APEX Calculus Textbook}
again at the end of the listed section.

Problems listed as PAULS-PRACT are found in \href{https://tutorial.math.lamar.edu/Problems/CalcI/CalcI.aspx}{Paul's Online Math Notes - Review Problems}
while problems listed as PAULS-ASSIGN are found in \href{https://tutorial.math.lamar.edu/ProblemsNS/CalcI/CalcI.aspx}{Paul's Online Math Notes - Assignment Problems}
and are indexed via [Topic]-[Section]-[Problem Number]

\section{Type 1 - Limits at Infinity}

Choose 4 of the following problems to solve.

\begin{enumerate}
  \item PAULS-ASSIGN.Limits.Limits at Infinity Part 1.4-17 (each of these counts as one problem)
  \item APEX 1.6.19-28 (each of these counts as one problem)
\end{enumerate}

\section{Type 2 - Limits Diverging to Infinity}

Choose 4 of the following problems to solve.

\begin{enumerate}
  \item PAULS-ASSIGN.Limits.Infinte Limits.1-8 (each of these counts as one problem)
  \item PAULS-ASSIGN.Limits.Infinte Limits.9-12 (each of these counts as one problem)
\end{enumerate}

\section{Type 3 - The Squeeze Theorem}

Choose 2 of the following problems to solve.

\begin{enumerate}
  \item APEX 1.3.35-38 (each of these counts as one problem)
  \item PAULS-ASSIGN.Limits.Computing Limits.31-35 (each of these counts as one problem)
\end{enumerate}

\section{Challenge Problem}
Recall the following facts/definitions:
\begin{enumerate}
  \item The \textit{degree} of a polynomial $p(x)$ is the highest power appearing in the polynomial, and it is denoted by $\mathrm{deg}(p(x))$. E.g. $\mathrm{deg}(x^2) = 2, \mathrm{deg}(x^5 + 4x^2 + 2x + 1) = 5$, etc.
  \item The \textit{leading term} of a polynomial $p(x)$ is the monomial summand with the highest degree, and it is denoted by $LT(p(x))$. E.g. $LT(x^3 + 3x^2 + 2x+1) = x^3, LT(10x^5 + 3x^2 + 1) = 10x^5$ etc.
  \item The \textit{leading coefficient} of a polynomial $p(x)$ is the coefficient of the leading term, and it is denoted by $LC(p(x))$.
  \item A \textit{rational function} is a function that can be expressed as the ratio of two polynomials.
\end{enumerate}

Show that the following facts about limits at infinity of a rational function 
\begin{equation*}
  \dfrac{p(x)}{q(x)} = \dfrac{a_n x^n + a_{n-1}x^{n-1} + \ldots + a_1 x + a_0}{b_m x^m + b_{m-1}x^{m-1} + \ldots + b_1 x + b_0} = \dfrac{\sum\limits_{i = 0}^n a_i x^i}{\sum\limits_{j = 0}^m b_j x^j}
\end{equation*}
hold:

\begin{enumerate}
  \item If $\mathrm{deg}(p(x)) > \mathrm{deg}(q(x))$ then:
  \begin{enumerate}

    \item If $\mathrm{deg}(p(x))$ is an \textit{even} number, 
      \begin{equation}
        \lim_{x \rightarrow \infty} \dfrac{p(x)}{q(x)} = \lim_{x \rightarrow -\infty} \dfrac{p(x)}{q(x)} = \begin{cases} \infty, & LC(p(x)) > 0 \\ -\infty, & LC(p(x)) < 0 \end{cases}
      \end{equation}

    \item If $\mathrm{deg}(p(x))$ is an \textit{odd} number,
      \begin{equation}
        \lim_{x \rightarrow \infty} \dfrac{p(x)}{q(x)} = - \lim_{x \rightarrow -\infty} \dfrac{p(x)}{q(x)} = \begin{cases} \infty, & LC(p(x)) > 0 \\ -\infty, & LC(p(x)) < 0 \end{cases}
      \end{equation}
  \end{enumerate}

  \item If $\mathrm{deg}(p(x)) = \mathrm{deg}(q(x))$ then
  \begin{equation}
    \lim_{x \rightarrow \infty} \dfrac{p(x)}{q(x)} = \lim_{x \rightarrow -\infty} \dfrac{p(x)}{q(x)} = \dfrac{LC(p(x))}{LC(q(x))}
  \end{equation}

  \item If $\mathrm{deg}(p(x)) < \mathrm{deg}(q(x))$ then
    \begin{equation}
      \lim_{x \rightarrow \infty} \dfrac{p(x)}{q(x)} = \lim_{x \rightarrow - \infty} \dfrac{p(x)}{q(x)} = 0
    \end{equation}
\end{enumerate}
% Most people don't use these, so they are "commented out"
% by starting the lines with a "%"
%\begin{keywords}
%   \LaTeX, typesetting
%\end{keywords}

%\begin{AMS}
%   50C60, 18C25
%\end{AMS}

%%%%%%%%%%%%%%%%%%%%%%
% % Here is the start of the Text
%%%%%%%%%%%%%%%%%%%%%%

\end{document}
